The indifference instrument forbids uncaused structure (distinctions nothing in the framework determines); it is a structural premise of the series, not a theorem, and is flagged as such wherever it is used. The identification of the three generations with the three two-twistor faces of the maximal composite is the families paper's Conjecture 3.1, at conjectural level, and the maximality it invokes is the Standard-Model paper's three-twistor bound, imported there and here at the observation level that remark declares; the mass bracket is face-valued, one bracket per face, so a vacuum is a direction in the three-dimensional face space, and its alignment is what breaks $S_{3}$ (structural level). Under the instrument the condensate's source lies on the democratic line spanned by $u = (1,1,1)/\sqrt{3}$, and, under the Gram-square reading of the leading mass form, the leading Yukawa form is $Y_{0} \propto \J = u'u'^{\mathsf{T}}$ with $u' = (1,1,1)$ and $\J$ the $3 \times 3$ all-ones matrix, of spectrum $(3, 0, 0)$: one heavy family and two light at leading order, lifted by the second-order breaking $S_{3} \to S_{2} \to 1$ that converts the CKM shape into CKM numbers. \emph{Provenance:} \cite[Fact~2.5, Conj.~3.1, Prop.~4.3, Prop.~4.4(1)]{ButlerT13}, each at the level stated there; \cite[Rem.~3.3]{ButlerT12} for the three-twistor bound, at observation level.
